\begin{lemma}{Norm und duale Norm}{NormUndDualeNorm}
Seien $(E, \norm\cdot_E)$ ein normierter $\K$--Vektorraum und $x \in E$.\\
Dann existiert ein $\vi \in \kran{E'}01$ 
mit $\norm x_E = \vi(x)$.\\
Insbesondere ist damit $\opnorm {ev(x, \cdot)} {E'} \K  = \norm x_E$.\\
Für alle $K \in \meng{\kabg{E'}01, \kran{E'}01}$
gilt daher $\norm x _ E = \sup\menge{\abs{\vi(x)}}{\vi \in K}$.\\
Für $\K = \R$ gilt außerdem $\norm x _ E = \sup\menge{\vi(x)}{\vi \in K}$.
\end{lemma}
\begin{proof}
Wir verwenden den folgenden Spezialfall des Satzes von Hahn--Banach:
\begin{quote}
Seien $(X, \norm\cdot)$ ein normierter $\K$--Vektorraum, $Y \le X$ und
$l \in Y'$.\\
Dann existiert ein $L \in X'$ mit $L\vert _Y = l$ und $\opnorm L X \K = \opnorm l Y \K$.
\end{quote}
Falls $x = 0$, so ist nichts zu tun, denn es gibt stetige Funktionale außer dem
Nullfunktional (wegen $E \neq 0$ kann man $y \in E\setminus \meng 0$ wählen
und damit wie unten beschrieben verfahren).\\
Sei also zusätzlich $x \neq 0$ und $x' \coloneqq  \frac 1 {\norm x _E} x$.\\
Seien $F\coloneqq \erz {x'}_\K$ und $\iota : \K \rightarrow F, k \mapsto k\cdot x'$.\\
Dann ist $\iota$ wegen $\norm{x'}_E=1$ eine bijektive, lineare Isometrie.\\
Sei nun $l \coloneqq  \iota\inv$. Dann ist $l$ ebenfalls eine lineare Isometrie, also
$\opnorm l F \K = 1$.\\
Also existiert nach obigem Satz ein $L \in E'$ mit $L \vert_ F = l$ und
$\opnorm L E \K = 1$.\\
Wegen $x \in F$ gilt weiter $L(x) = L(\norm x_E \cdot x') = \norm x_E$, womit
der erste Teil der Behauptung gezeigt ist.\\
Seien $K \in \meng{\kabg{E'}01, \kran{E'}01}$ und $\vi \in K$.\\
Dann gilt $\abs{\vi(x)} \le \opnorm \vi E \K \cdot \norm x_E \le \norm x_E$.\\
Damit ist bereits $\opnorm {ev(x, \cdot)} {E'} \K \le \norm x _E$.\\
Da aber weiter $L \in K$ ist, gilt auch Gleichheit, denn nach Definition und
einer aus dem Grundstudium bekannten Rechnung ist
\[\opnorm {ev(x, \cdot)} {E'} \K = \sup\menge{\abs{ev(x,\psi)}}{\psi \in K} =
                          \sup\menge{\abs{\psi(x)}}{\psi \in K}\text.\]
Im Falle $\K = \R$ kann man offensichtlich auf den Betrag verzichten.
\end{proof}